\section{Testes, ferramentas e geração das figuras}
\label{app:repro}

Os testes automatizados ficam em \code{tests/} e são compilados pelo CMake; as ferramentas do relatório ficam em
\code{doc/report/tools/}. Todos usam o código do repositório como está e podem ser rodados a partir da raiz do
repositório. Cada arquivo começa com um comentário que descreve seu uso; este apêndice explica o que cada um faz e como
os resultados do relatório foram obtidos.

\begin{center}\small
\begin{tabularx}{\linewidth}{>{\ttfamily\small}p{4.6cm}X}
\toprule
\normalfont Arquivo & Papel\\
\midrule
tests/test\_*.cpp & suítes do Google Test (\cref{tab:suites});\\
tests/support/ & malhas de teste, verificações V1--V7 e acesso de leitura ao estado do construtor, usados pelos testes e
pelo \code{driver};\\
tests/CMakeLists.txt & registra as suítes e o teste com malhas aleatórias no CTest;\\
tools/driver.cpp & programa de teste da \code{IBHMInterfaceElementBuilder}: insere arestas, verifica a malha e grava cada
estado em JSON;\\
tools/compare.py & compara a topologia final de execuções diferentes;\\
tools/stress.py & teste com malhas aleatórias, várias ordens e inserção em lote;\\
tools/gen.py & gera as figuras de malha do relatório (\code{doc/report/figs/}).\\
\bottomrule
\end{tabularx}
\end{center}

\subsection{Como compilar e rodar}

Os comandos abaixo partem da raiz do repositório; \code{\$W} é um diretório de trabalho qualquer, fora do repositório.
O CMake não depende de Qt. O Google Test vem do pacote instalado, se houver; senão, a versão 1.17.0 é baixada na
configuração, com o hash SHA256 fixado no \code{CMakeLists.txt}. (No macOS, o \code{cmake} e o \code{ninja} do CLion
ficam em \code{/Applications/CLion.app/Contents/bin/}.) O tipo de build pode ser \code{Debug} ou \code{Release}: nos
dois, os testes rodam com as asserções ligadas.
\begin{lstlisting}[language=bash]
cmake -S . -B $W/build -G Ninja -DCMAKE_BUILD_TYPE=Debug
cmake --build $W/build
cd $W/build && ctest -j8                      # ou: ctest -L unit; ctest -L stress
D=$W/build/tests/driver
$D malha2Hyb.txt $W/ex.json 31 34 11 26 47 38 23 19 30 42 29 45         # exemplo dos slides
python3 doc/report/tools/stress.py $D 21 400 $W/s1                       # rodada 1
python3 doc/report/tools/stress.py $D 22 400 $W/s2 --skip-border-multi   # rodada 2
python3 doc/report/tools/gen.py $D doc/report/figs $W/work               # figuras
cd doc/report && tectonic -X compile relatorio.tex                       # relatório
\end{lstlisting}
Sem CMake, o \code{driver} também pode ser compilado diretamente:
\begin{lstlisting}[language=bash]
c++ -std=c++17 -Icore/TopologyDataStructures -Icore/InterfaceElementsBuilder -Itests/support \
    doc/report/tools/driver.cpp tests/support/MeshValidation.cpp \
    core/TopologyDataStructures/IBHM.cpp \
    core/InterfaceElementsBuilder/IBHMInterfaceElementBuilder.cpp -o $W/driver
\end{lstlisting}

\subsection{Os testes automatizados (\code{tests/})}

O alvo \code{InterfaceElementsTests} reúne as suítes da \cref{tab:suites}, e \code{gtest\_discover\_tests} registra
cada \code{TEST} como um teste do CTest, com o rótulo \code{unit} e limite de 60\,s. Como cada teste roda num processo
próprio, uma asserção ou um laço infinito faz falhar apenas aquele teste. O teste com malhas aleatórias entra como
\code{stress\_supported} (rótulo \code{stress}), que chama o \code{stress.py} com o \code{driver}; o
\code{stress\_all} está registrado, mas desativado até o I1 ser corrigido.

\textbf{Asserções ligadas.} Alguns defeitos só aparecem como asserção, e o teste de morte de \code{addPoint} espera uma.
Por isso o \code{tests/CMakeLists.txt} retira o \code{NDEBUG} das opções do diretório e compila o núcleo de novo
(\code{InterfaceElementsLibChecked}), qualquer que seja o tipo de build; a biblioteca \code{InterfaceElementsLib}
continua compilada normalmente.

\textbf{Código de apoio} (\code{tests/support/}):
\begin{itemize}
\item \code{TestMeshes}: a malha base (\code{malha2Hyb.txt}, com o vértice $T$ em $(4,3)$) e a lista das 12 arestas do
exemplo, a malha em leque do caso de borda com duas arestas e uma malha de dois triângulos;
\item \code{MeshValidation}: a ``foto'' da malha (lida pela interface pública da \code{IBHM}, com $\nxt$/$\prv$
recalculados), as verificações V1--V5, a origem de cada cópia e a V7, e a partição dos cantos usada para comparar
topologias;
\item \code{BuilderInspector}: acesso de leitura ao estado privado do construtor (o operador que seria escolhido,
\code{\_isInterfaceElement}, \code{\_inInterfaceElement} e \code{\_collapsed}). O construtor a declara como classe
amiga. A alternativa usada até a v0.3, \code{\#define private public}, foi abandonada: no MSVC o nível de acesso faz
parte do nome decorado das funções, e a ligação falha.
\end{itemize}

\textbf{Defeitos em aberto.} Os testes de \code{test\_known\_defects.cpp} têm o prefixo \code{DISABLED\_}: aparecem como
desativados em toda execução e descrevem o comportamento correto, de modo que, corrigido o defeito, basta tirar o
prefixo. Para rodá-los mesmo assim: \code{InterfaceElementsTests --gtest\_also\_run\_disabled\_tests
--gtest\_filter='KnownDefects*'}.

\subsection{O programa de teste (\code{driver.cpp})}

\textbf{Uso:} \code{driver [--batch] <malha> <saida.json> <he0> <he1> ...}. A malha está no formato de
\code{malha2Hyb.txt}: o número de vértices e de elementos, as coordenadas e, para cada elemento, o número de vértices
seguido dos índices. Cada \code{he} é uma half-edge da aresta a abrir (qualquer uma das duas). Sem \code{--batch}, cada
aresta é inserida numa chamada de \code{insertInterfaceElements}; com \code{--batch}, todas vão numa única chamada.

\textbf{O que faz em cada chamada:}
\begin{enumerate}
\item tira uma ``foto'' da malha (vértices de cada elemento, offsets e opostas), lida pela interface pública da
\code{IBHM};
\item se a aresta vai ser aceita, pergunta a \code{retrieveOperator} qual operador será usado em cada extremidade, só para
registrá-lo (a decisão não muda nada na malha);
\item chama \code{computeNumberOfNewVertices} e \code{insertInterfaceElements};
\item descobre o vértice de origem de cada vértice novo: todo vértice novo substitui, em algum canto de um elemento que
já existia, o vértice que estava lá antes; esse vértice antigo (ou a origem dele) é a origem do novo;
\item aplica as verificações V1--V5 e V7 da \cref{sec:verificacao} e compara o número de vértices criados com o
previsto (V6);
\item imprime uma linha e grava o estado no JSON.
\end{enumerate}
As verificações são feitas sobre a ``foto'', com consultas $\nxt$/$\prv$ recalculadas, para não depender do código que
está sendo testado; a foto, as verificações e a origem das cópias vêm de \code{tests/support/MeshValidation}, as mesmas
dos testes automatizados. O estado privado do construtor (\code{\_collapsed}, \code{\_isInterfaceElement} e o operador
escolhido) é lido por \code{BuilderInspector}. Na v0.4 o \code{driver} passou a usar esse código compartilhado; sua saída
(console e JSON) continuou idêntica, byte a byte, no exemplo, nos casos de entrada inválida e em 834 execuções do
estresse (120 malhas de cada rodada, inclusive as que abortam).

\textbf{Saída no console:} uma linha por chamada, como
\begin{lstlisting}[language={}]
step  2 edge he=34 (10-11)  A:SPLIT_ELEMENT B:CANONICAL     new=1 predicted=1 ok
\end{lstlisting}
com a half-edge, os dois vértices da aresta, o operador de cada extremidade, os vértices criados e previstos e o
resultado (\code{ok} ou \code{INVALID}, seguido das verificações que falharam). O código de saída é 0 se tudo passou e
2 se alguma verificação falhou; uma asserção da biblioteca interrompe o programa.

\textbf{JSON:} um objeto com \code{points} (coordenadas originais) e \code{steps}, uma lista em que o passo 0 é a malha de
entrada e o passo $k$ é o estado depois da $k$-ésima chamada. Cada passo tem:

\begin{center}\small
\begin{tabularx}{\linewidth}{>{\ttfamily\small}p{3.6cm}X}
\toprule
\normalfont Campo & Conteúdo\\
\midrule
edge, A, B & a half-edge inserida e os vértices da sua aresta\\
opA, opB & o operador de cada extremidade (vazio no modo em lote ou para arestas rejeitadas)\\
predicted, created & vértices previstos por \code{computeNumberOfNewVertices} e vértices criados\\
errors, msgs & número de problemas e mensagens das verificações\\
elements, first & os vértices de cada elemento e a primeira half-edge de cada elemento (os offsets)\\
opp & a oposta de cada half-edge; $-1$ é \BORDER{} e $-2$ é \COLLAPSED\\
isIE & 1 para os elementos de interface\\
root & para cada vértice, o vértice original de que ele é cópia\\
collapsed & o mapa \code{\_collapsed} do construtor, como pares (vértice, half-edge)\\
\bottomrule
\end{tabularx}
\end{center}
O arquivo é gravado e esvaziado (\emph{flush}) a cada passo; se um passo posterior aborta, os anteriores continuam
legíveis (o \code{gen.py} usa isso para desenhar o caso de borda que falha).

\subsection{Comparação de topologias (\code{compare.py})}

Duas execuções que inserem as mesmas arestas em ordens diferentes criam os vértices novos em ordens diferentes, então os
índices não servem para comparar. O \code{compare.py} olha só para os elementos que já existiam na malha de entrada
(eles mantêm seus índices) e agrupa os cantos desses elementos pelo vértice que usam. O conjunto de grupos é uma partição
dos cantos; duas execuções têm a mesma topologia, a menos de renumeração, se e somente se as partições são iguais. Os
elementos de interface ficam determinados pelos elementos contínuos em volta e não precisam entrar na comparação. Pela
linha de comando, \code{python3 compare.py a.json b.json ...} diz, para cada arquivo, se a partição é igual à do
primeiro.

\subsection{Teste com malhas aleatórias (\code{stress.py})}

\textbf{Uso:} \code{stress.py <driver> <semente> <casos> <diretório> [--skip-border-multi]}. Cada caso:
\begin{enumerate}
\item gera uma malha numa grade de $3$ a $7$ por $3$ a $7$ células, alternando quatro tipos (\cref{fig:estresse-malhas}):
só quadriláteros, só triângulos (cada célula dividida por uma das diagonais, ao acaso), híbrida (60\% das células viram
triângulos) e com polígonos (como a híbrida, mas 25\% das vezes duas células vizinhas viram um hexágono);
\item sorteia até metade das arestas internas;
\item insere o mesmo conjunto quatro vezes com o \code{driver}: três vezes uma aresta por chamada, em ordens aleatórias e
com a half-edge de cada aresta escolhida ao acaso, e uma vez em lote;
\item exige que as quatro execuções sejam válidas, que a contagem de vértices seja sempre a prevista e que as quatro
partições sejam iguais.
\end{enumerate}
Um caso é marcado como ``vértice de borda com 2+ arestas'' quando a seleção tem um vértice de borda tocado por duas ou
mais arestas (a configuração do defeito I1). Com \code{--skip-border-multi}, essas seleções são sorteadas de novo. No
fim, o script imprime os contadores (casos válidos, inválidos e abortados, quantos abortados tinham a configuração do I1,
casos dependentes da ordem, casos em que o lote difere, contagens erradas, operadores aplicados), a asserção que causou
cada aborto e o número de casos válidos por tipo de malha. O código de saída é 0 se todos os casos passaram e 1 se
houve aborto, caso inválido, dependência da ordem ou do lote, ou contagem errada; é ele que o CTest usa. A semente fixa
o gerador aleatório: a mesma semente sempre gera as mesmas malhas, seleções e ordens (o relatório usa as sementes 21 e
22, com 400 casos cada).

\begin{figure}[ht]
\centering
\begin{subfigure}[b]{0.24\linewidth}\centering\resizebox{\linewidth}{!}{\begin{tikzpicture}[scale=0.620, every node/.style={font=\scriptsize}]
\filldraw[cefill] (0.000,0.000) -- (1.000,0.000) -- (1.000,1.000) -- (0.000,1.000) -- cycle;
\filldraw[cefill] (1.000,0.000) -- (2.000,0.000) -- (2.000,1.000) -- (1.000,1.000) -- cycle;
\filldraw[cefill] (2.000,0.000) -- (3.000,0.000) -- (3.000,1.000) -- (2.000,1.000) -- cycle;
\filldraw[cefill] (3.000,0.000) -- (4.000,0.000) -- (4.000,1.000) -- (3.000,1.000) -- cycle;
\filldraw[cefill] (4.000,0.000) -- (5.000,0.000) -- (5.000,1.000) -- (4.000,1.000) -- cycle;
\filldraw[cefill] (5.000,0.000) -- (6.000,0.000) -- (6.000,1.000) -- (5.000,1.000) -- cycle;
\filldraw[cefill] (0.000,1.000) -- (1.000,1.000) -- (1.000,2.000) -- (0.000,2.000) -- cycle;
\filldraw[cefill] (1.000,1.000) -- (2.000,1.000) -- (2.000,2.000) -- (1.000,2.000) -- cycle;
\filldraw[cefill] (2.000,1.000) -- (3.000,1.000) -- (3.000,2.000) -- (2.000,2.000) -- cycle;
\filldraw[cefill] (3.000,1.000) -- (4.000,1.000) -- (4.000,2.000) -- (3.000,2.000) -- cycle;
\filldraw[cefill] (4.000,1.000) -- (5.000,1.000) -- (5.000,2.000) -- (4.000,2.000) -- cycle;
\filldraw[cefill] (5.000,1.000) -- (6.000,1.000) -- (6.000,2.000) -- (5.000,2.000) -- cycle;
\filldraw[cefill] (0.000,2.000) -- (1.000,2.000) -- (1.000,3.000) -- (0.000,3.000) -- cycle;
\filldraw[cefill] (1.000,2.000) -- (2.000,2.000) -- (2.000,3.000) -- (1.000,3.000) -- cycle;
\filldraw[cefill] (2.000,2.000) -- (3.000,2.000) -- (3.000,3.000) -- (2.000,3.000) -- cycle;
\filldraw[cefill] (3.000,2.000) -- (4.000,2.000) -- (4.000,3.000) -- (3.000,3.000) -- cycle;
\filldraw[cefill] (4.000,2.000) -- (5.000,2.000) -- (5.000,3.000) -- (4.000,3.000) -- cycle;
\filldraw[cefill] (5.000,2.000) -- (6.000,2.000) -- (6.000,3.000) -- (5.000,3.000) -- cycle;
\filldraw[cefill] (0.000,3.000) -- (1.000,3.000) -- (1.000,4.000) -- (0.000,4.000) -- cycle;
\filldraw[cefill] (1.000,3.000) -- (2.000,3.000) -- (2.000,4.000) -- (1.000,4.000) -- cycle;
\filldraw[cefill] (2.000,3.000) -- (3.000,3.000) -- (3.000,4.000) -- (2.000,4.000) -- cycle;
\filldraw[cefill] (3.000,3.000) -- (4.000,3.000) -- (4.000,4.000) -- (3.000,4.000) -- cycle;
\filldraw[cefill] (4.000,3.000) -- (5.000,3.000) -- (5.000,4.000) -- (4.000,4.000) -- cycle;
\filldraw[cefill] (5.000,3.000) -- (6.000,3.000) -- (6.000,4.000) -- (5.000,4.000) -- cycle;
\draw[seledge] (2.000,2.000) -- (2.000,3.000);
\draw[seledge] (2.000,0.000) -- (2.000,1.000);
\draw[seledge] (1.000,2.000) -- (1.000,3.000);
\draw[seledge] (2.000,3.000) -- (2.000,4.000);
\draw[seledge] (4.000,2.000) -- (5.000,2.000);
\draw[seledge] (0.000,3.000) -- (1.000,3.000);
\draw[seledge] (5.000,2.000) -- (5.000,3.000);
\draw[seledge] (4.000,2.000) -- (4.000,3.000);
\draw[seledge] (2.000,3.000) -- (3.000,3.000);
\node[vreg, dot] at (0.000,0.000) {};
\node[vreg, dot] at (1.000,0.000) {};
\node[vreg, dot] at (2.000,0.000) {};
\node[vreg, dot] at (3.000,0.000) {};
\node[vreg, dot] at (4.000,0.000) {};
\node[vreg, dot] at (5.000,0.000) {};
\node[vreg, dot] at (6.000,0.000) {};
\node[vreg, dot] at (0.000,1.000) {};
\node[vreg, dot] at (1.000,1.000) {};
\node[vreg, dot] at (2.000,1.000) {};
\node[vreg, dot] at (3.000,1.000) {};
\node[vreg, dot] at (4.000,1.000) {};
\node[vreg, dot] at (5.000,1.000) {};
\node[vreg, dot] at (6.000,1.000) {};
\node[vreg, dot] at (0.000,2.000) {};
\node[vreg, dot] at (1.000,2.000) {};
\node[vreg, dot] at (2.000,2.000) {};
\node[vreg, dot] at (3.000,2.000) {};
\node[vreg, dot] at (4.000,2.000) {};
\node[vreg, dot] at (5.000,2.000) {};
\node[vreg, dot] at (6.000,2.000) {};
\node[vreg, dot] at (0.000,3.000) {};
\node[vreg, dot] at (1.000,3.000) {};
\node[vreg, dot] at (2.000,3.000) {};
\node[vreg, dot] at (3.000,3.000) {};
\node[vreg, dot] at (4.000,3.000) {};
\node[vreg, dot] at (5.000,3.000) {};
\node[vreg, dot] at (6.000,3.000) {};
\node[vreg, dot] at (0.000,4.000) {};
\node[vreg, dot] at (1.000,4.000) {};
\node[vreg, dot] at (2.000,4.000) {};
\node[vreg, dot] at (3.000,4.000) {};
\node[vreg, dot] at (4.000,4.000) {};
\node[vreg, dot] at (5.000,4.000) {};
\node[vreg, dot] at (6.000,4.000) {};

\end{tikzpicture}
}\caption{quad}\end{subfigure}\hfill
\begin{subfigure}[b]{0.24\linewidth}\centering\resizebox{\linewidth}{!}{\begin{tikzpicture}[scale=0.620, every node/.style={font=\scriptsize}]
\filldraw[cefill] (0.000,0.000) -- (1.000,0.000) -- (0.000,1.000) -- cycle;
\filldraw[cefill] (1.000,0.000) -- (1.000,1.000) -- (0.000,1.000) -- cycle;
\filldraw[cefill] (1.000,0.000) -- (2.000,0.000) -- (1.000,1.000) -- cycle;
\filldraw[cefill] (2.000,0.000) -- (2.000,1.000) -- (1.000,1.000) -- cycle;
\filldraw[cefill] (2.000,0.000) -- (3.000,0.000) -- (3.000,1.000) -- cycle;
\filldraw[cefill] (2.000,0.000) -- (3.000,1.000) -- (2.000,1.000) -- cycle;
\filldraw[cefill] (3.000,0.000) -- (4.000,0.000) -- (3.000,1.000) -- cycle;
\filldraw[cefill] (4.000,0.000) -- (4.000,1.000) -- (3.000,1.000) -- cycle;
\filldraw[cefill] (4.000,0.000) -- (5.000,0.000) -- (5.000,1.000) -- cycle;
\filldraw[cefill] (4.000,0.000) -- (5.000,1.000) -- (4.000,1.000) -- cycle;
\filldraw[cefill] (5.000,0.000) -- (6.000,0.000) -- (6.000,1.000) -- cycle;
\filldraw[cefill] (5.000,0.000) -- (6.000,1.000) -- (5.000,1.000) -- cycle;
\filldraw[cefill] (0.000,1.000) -- (1.000,1.000) -- (1.000,2.000) -- cycle;
\filldraw[cefill] (0.000,1.000) -- (1.000,2.000) -- (0.000,2.000) -- cycle;
\filldraw[cefill] (1.000,1.000) -- (2.000,1.000) -- (2.000,2.000) -- cycle;
\filldraw[cefill] (1.000,1.000) -- (2.000,2.000) -- (1.000,2.000) -- cycle;
\filldraw[cefill] (2.000,1.000) -- (3.000,1.000) -- (2.000,2.000) -- cycle;
\filldraw[cefill] (3.000,1.000) -- (3.000,2.000) -- (2.000,2.000) -- cycle;
\filldraw[cefill] (3.000,1.000) -- (4.000,1.000) -- (3.000,2.000) -- cycle;
\filldraw[cefill] (4.000,1.000) -- (4.000,2.000) -- (3.000,2.000) -- cycle;
\filldraw[cefill] (4.000,1.000) -- (5.000,1.000) -- (5.000,2.000) -- cycle;
\filldraw[cefill] (4.000,1.000) -- (5.000,2.000) -- (4.000,2.000) -- cycle;
\filldraw[cefill] (5.000,1.000) -- (6.000,1.000) -- (5.000,2.000) -- cycle;
\filldraw[cefill] (6.000,1.000) -- (6.000,2.000) -- (5.000,2.000) -- cycle;
\filldraw[cefill] (0.000,2.000) -- (1.000,2.000) -- (1.000,3.000) -- cycle;
\filldraw[cefill] (0.000,2.000) -- (1.000,3.000) -- (0.000,3.000) -- cycle;
\filldraw[cefill] (1.000,2.000) -- (2.000,2.000) -- (2.000,3.000) -- cycle;
\filldraw[cefill] (1.000,2.000) -- (2.000,3.000) -- (1.000,3.000) -- cycle;
\filldraw[cefill] (2.000,2.000) -- (3.000,2.000) -- (2.000,3.000) -- cycle;
\filldraw[cefill] (3.000,2.000) -- (3.000,3.000) -- (2.000,3.000) -- cycle;
\filldraw[cefill] (3.000,2.000) -- (4.000,2.000) -- (3.000,3.000) -- cycle;
\filldraw[cefill] (4.000,2.000) -- (4.000,3.000) -- (3.000,3.000) -- cycle;
\filldraw[cefill] (4.000,2.000) -- (5.000,2.000) -- (4.000,3.000) -- cycle;
\filldraw[cefill] (5.000,2.000) -- (5.000,3.000) -- (4.000,3.000) -- cycle;
\filldraw[cefill] (5.000,2.000) -- (6.000,2.000) -- (5.000,3.000) -- cycle;
\filldraw[cefill] (6.000,2.000) -- (6.000,3.000) -- (5.000,3.000) -- cycle;
\filldraw[cefill] (0.000,3.000) -- (1.000,3.000) -- (1.000,4.000) -- cycle;
\filldraw[cefill] (0.000,3.000) -- (1.000,4.000) -- (0.000,4.000) -- cycle;
\filldraw[cefill] (1.000,3.000) -- (2.000,3.000) -- (2.000,4.000) -- cycle;
\filldraw[cefill] (1.000,3.000) -- (2.000,4.000) -- (1.000,4.000) -- cycle;
\filldraw[cefill] (2.000,3.000) -- (3.000,3.000) -- (3.000,4.000) -- cycle;
\filldraw[cefill] (2.000,3.000) -- (3.000,4.000) -- (2.000,4.000) -- cycle;
\filldraw[cefill] (3.000,3.000) -- (4.000,3.000) -- (4.000,4.000) -- cycle;
\filldraw[cefill] (3.000,3.000) -- (4.000,4.000) -- (3.000,4.000) -- cycle;
\filldraw[cefill] (4.000,3.000) -- (5.000,3.000) -- (5.000,4.000) -- cycle;
\filldraw[cefill] (4.000,3.000) -- (5.000,4.000) -- (4.000,4.000) -- cycle;
\filldraw[cefill] (5.000,3.000) -- (6.000,3.000) -- (5.000,4.000) -- cycle;
\filldraw[cefill] (6.000,3.000) -- (6.000,4.000) -- (5.000,4.000) -- cycle;
\draw[seledge] (3.000,2.000) -- (4.000,2.000);
\draw[seledge] (1.000,1.000) -- (2.000,1.000);
\draw[seledge] (4.000,3.000) -- (5.000,3.000);
\draw[seledge] (4.000,1.000) -- (3.000,2.000);
\draw[seledge] (5.000,2.000) -- (6.000,2.000);
\draw[seledge] (1.000,0.000) -- (1.000,1.000);
\draw[seledge] (3.000,3.000) -- (4.000,3.000);
\draw[seledge] (2.000,1.000) -- (3.000,1.000);
\draw[seledge] (4.000,2.000) -- (3.000,3.000);
\draw[seledge] (5.000,2.000) -- (5.000,3.000);
\draw[seledge] (6.000,2.000) -- (5.000,3.000);
\draw[seledge] (5.000,0.000) -- (6.000,1.000);
\draw[seledge] (2.000,3.000) -- (3.000,3.000);
\draw[seledge] (3.000,1.000) -- (4.000,1.000);
\draw[seledge] (2.000,2.000) -- (3.000,2.000);
\node[vreg, dot] at (0.000,0.000) {};
\node[vreg, dot] at (1.000,0.000) {};
\node[vreg, dot] at (2.000,0.000) {};
\node[vreg, dot] at (3.000,0.000) {};
\node[vreg, dot] at (4.000,0.000) {};
\node[vreg, dot] at (5.000,0.000) {};
\node[vreg, dot] at (6.000,0.000) {};
\node[vreg, dot] at (0.000,1.000) {};
\node[vreg, dot] at (1.000,1.000) {};
\node[vreg, dot] at (2.000,1.000) {};
\node[vreg, dot] at (3.000,1.000) {};
\node[vreg, dot] at (4.000,1.000) {};
\node[vreg, dot] at (5.000,1.000) {};
\node[vreg, dot] at (6.000,1.000) {};
\node[vreg, dot] at (0.000,2.000) {};
\node[vreg, dot] at (1.000,2.000) {};
\node[vreg, dot] at (2.000,2.000) {};
\node[vreg, dot] at (3.000,2.000) {};
\node[vreg, dot] at (4.000,2.000) {};
\node[vreg, dot] at (5.000,2.000) {};
\node[vreg, dot] at (6.000,2.000) {};
\node[vreg, dot] at (0.000,3.000) {};
\node[vreg, dot] at (1.000,3.000) {};
\node[vreg, dot] at (2.000,3.000) {};
\node[vreg, dot] at (3.000,3.000) {};
\node[vreg, dot] at (4.000,3.000) {};
\node[vreg, dot] at (5.000,3.000) {};
\node[vreg, dot] at (6.000,3.000) {};
\node[vreg, dot] at (0.000,4.000) {};
\node[vreg, dot] at (1.000,4.000) {};
\node[vreg, dot] at (2.000,4.000) {};
\node[vreg, dot] at (3.000,4.000) {};
\node[vreg, dot] at (4.000,4.000) {};
\node[vreg, dot] at (5.000,4.000) {};
\node[vreg, dot] at (6.000,4.000) {};

\end{tikzpicture}
}\caption{tri}\end{subfigure}\hfill
\begin{subfigure}[b]{0.24\linewidth}\centering\resizebox{\linewidth}{!}{\begin{tikzpicture}[scale=0.620, every node/.style={font=\scriptsize}]
\filldraw[cefill] (0.000,0.000) -- (1.000,0.000) -- (0.000,1.000) -- cycle;
\filldraw[cefill] (1.000,0.000) -- (1.000,1.000) -- (0.000,1.000) -- cycle;
\filldraw[cefill] (1.000,0.000) -- (2.000,0.000) -- (1.000,1.000) -- cycle;
\filldraw[cefill] (2.000,0.000) -- (2.000,1.000) -- (1.000,1.000) -- cycle;
\filldraw[cefill] (2.000,0.000) -- (3.000,0.000) -- (3.000,1.000) -- (2.000,1.000) -- cycle;
\filldraw[cefill] (3.000,0.000) -- (4.000,0.000) -- (4.000,1.000) -- cycle;
\filldraw[cefill] (3.000,0.000) -- (4.000,1.000) -- (3.000,1.000) -- cycle;
\filldraw[cefill] (4.000,0.000) -- (5.000,0.000) -- (5.000,1.000) -- (4.000,1.000) -- cycle;
\filldraw[cefill] (5.000,0.000) -- (6.000,0.000) -- (6.000,1.000) -- cycle;
\filldraw[cefill] (5.000,0.000) -- (6.000,1.000) -- (5.000,1.000) -- cycle;
\filldraw[cefill] (0.000,1.000) -- (1.000,1.000) -- (1.000,2.000) -- (0.000,2.000) -- cycle;
\filldraw[cefill] (1.000,1.000) -- (2.000,1.000) -- (1.000,2.000) -- cycle;
\filldraw[cefill] (2.000,1.000) -- (2.000,2.000) -- (1.000,2.000) -- cycle;
\filldraw[cefill] (2.000,1.000) -- (3.000,1.000) -- (2.000,2.000) -- cycle;
\filldraw[cefill] (3.000,1.000) -- (3.000,2.000) -- (2.000,2.000) -- cycle;
\filldraw[cefill] (3.000,1.000) -- (4.000,1.000) -- (3.000,2.000) -- cycle;
\filldraw[cefill] (4.000,1.000) -- (4.000,2.000) -- (3.000,2.000) -- cycle;
\filldraw[cefill] (4.000,1.000) -- (5.000,1.000) -- (5.000,2.000) -- (4.000,2.000) -- cycle;
\filldraw[cefill] (5.000,1.000) -- (6.000,1.000) -- (5.000,2.000) -- cycle;
\filldraw[cefill] (6.000,1.000) -- (6.000,2.000) -- (5.000,2.000) -- cycle;
\filldraw[cefill] (0.000,2.000) -- (1.000,2.000) -- (1.000,3.000) -- (0.000,3.000) -- cycle;
\filldraw[cefill] (1.000,2.000) -- (2.000,2.000) -- (1.000,3.000) -- cycle;
\filldraw[cefill] (2.000,2.000) -- (2.000,3.000) -- (1.000,3.000) -- cycle;
\filldraw[cefill] (2.000,2.000) -- (3.000,2.000) -- (3.000,3.000) -- cycle;
\filldraw[cefill] (2.000,2.000) -- (3.000,3.000) -- (2.000,3.000) -- cycle;
\filldraw[cefill] (3.000,2.000) -- (4.000,2.000) -- (3.000,3.000) -- cycle;
\filldraw[cefill] (4.000,2.000) -- (4.000,3.000) -- (3.000,3.000) -- cycle;
\filldraw[cefill] (4.000,2.000) -- (5.000,2.000) -- (4.000,3.000) -- cycle;
\filldraw[cefill] (5.000,2.000) -- (5.000,3.000) -- (4.000,3.000) -- cycle;
\filldraw[cefill] (5.000,2.000) -- (6.000,2.000) -- (6.000,3.000) -- (5.000,3.000) -- cycle;
\filldraw[cefill] (0.000,3.000) -- (1.000,3.000) -- (1.000,4.000) -- (0.000,4.000) -- cycle;
\filldraw[cefill] (1.000,3.000) -- (2.000,3.000) -- (2.000,4.000) -- (1.000,4.000) -- cycle;
\filldraw[cefill] (2.000,3.000) -- (3.000,3.000) -- (2.000,4.000) -- cycle;
\filldraw[cefill] (3.000,3.000) -- (3.000,4.000) -- (2.000,4.000) -- cycle;
\filldraw[cefill] (3.000,3.000) -- (4.000,3.000) -- (3.000,4.000) -- cycle;
\filldraw[cefill] (4.000,3.000) -- (4.000,4.000) -- (3.000,4.000) -- cycle;
\filldraw[cefill] (4.000,3.000) -- (5.000,3.000) -- (5.000,4.000) -- (4.000,4.000) -- cycle;
\filldraw[cefill] (5.000,3.000) -- (6.000,3.000) -- (6.000,4.000) -- cycle;
\filldraw[cefill] (5.000,3.000) -- (6.000,4.000) -- (5.000,4.000) -- cycle;
\draw[seledge] (2.000,1.000) -- (2.000,2.000);
\draw[seledge] (3.000,1.000) -- (3.000,2.000);
\draw[seledge] (2.000,3.000) -- (3.000,3.000);
\draw[seledge] (2.000,2.000) -- (1.000,3.000);
\draw[seledge] (1.000,3.000) -- (1.000,4.000);
\draw[seledge] (4.000,1.000) -- (4.000,2.000);
\draw[seledge] (1.000,2.000) -- (2.000,2.000);
\draw[seledge] (3.000,2.000) -- (4.000,2.000);
\draw[seledge] (0.000,2.000) -- (1.000,2.000);
\draw[seledge] (5.000,2.000) -- (4.000,3.000);
\draw[seledge] (5.000,1.000) -- (6.000,1.000);
\draw[seledge] (4.000,2.000) -- (4.000,3.000);
\draw[seledge] (5.000,2.000) -- (5.000,3.000);
\node[vreg, dot] at (0.000,0.000) {};
\node[vreg, dot] at (1.000,0.000) {};
\node[vreg, dot] at (2.000,0.000) {};
\node[vreg, dot] at (3.000,0.000) {};
\node[vreg, dot] at (4.000,0.000) {};
\node[vreg, dot] at (5.000,0.000) {};
\node[vreg, dot] at (6.000,0.000) {};
\node[vreg, dot] at (0.000,1.000) {};
\node[vreg, dot] at (1.000,1.000) {};
\node[vreg, dot] at (2.000,1.000) {};
\node[vreg, dot] at (3.000,1.000) {};
\node[vreg, dot] at (4.000,1.000) {};
\node[vreg, dot] at (5.000,1.000) {};
\node[vreg, dot] at (6.000,1.000) {};
\node[vreg, dot] at (0.000,2.000) {};
\node[vreg, dot] at (1.000,2.000) {};
\node[vreg, dot] at (2.000,2.000) {};
\node[vreg, dot] at (3.000,2.000) {};
\node[vreg, dot] at (4.000,2.000) {};
\node[vreg, dot] at (5.000,2.000) {};
\node[vreg, dot] at (6.000,2.000) {};
\node[vreg, dot] at (0.000,3.000) {};
\node[vreg, dot] at (1.000,3.000) {};
\node[vreg, dot] at (2.000,3.000) {};
\node[vreg, dot] at (3.000,3.000) {};
\node[vreg, dot] at (4.000,3.000) {};
\node[vreg, dot] at (5.000,3.000) {};
\node[vreg, dot] at (6.000,3.000) {};
\node[vreg, dot] at (0.000,4.000) {};
\node[vreg, dot] at (1.000,4.000) {};
\node[vreg, dot] at (2.000,4.000) {};
\node[vreg, dot] at (3.000,4.000) {};
\node[vreg, dot] at (4.000,4.000) {};
\node[vreg, dot] at (5.000,4.000) {};
\node[vreg, dot] at (6.000,4.000) {};

\end{tikzpicture}
}\caption{hybrid}\end{subfigure}\hfill
\begin{subfigure}[b]{0.24\linewidth}\centering\resizebox{\linewidth}{!}{\begin{tikzpicture}[scale=0.620, every node/.style={font=\scriptsize}]
\filldraw[cefill] (0.000,0.000) -- (1.000,0.000) -- (2.000,0.000) -- (2.000,1.000) -- (1.000,1.000) -- (0.000,1.000) -- cycle;
\filldraw[cefill] (2.000,0.000) -- (3.000,0.000) -- (3.000,1.000) -- cycle;
\filldraw[cefill] (2.000,0.000) -- (3.000,1.000) -- (2.000,1.000) -- cycle;
\filldraw[cefill] (3.000,0.000) -- (4.000,0.000) -- (4.000,1.000) -- (3.000,1.000) -- cycle;
\filldraw[cefill] (4.000,0.000) -- (5.000,0.000) -- (5.000,1.000) -- (4.000,1.000) -- cycle;
\filldraw[cefill] (5.000,0.000) -- (6.000,0.000) -- (6.000,1.000) -- cycle;
\filldraw[cefill] (5.000,0.000) -- (6.000,1.000) -- (5.000,1.000) -- cycle;
\filldraw[cefill] (0.000,1.000) -- (1.000,1.000) -- (1.000,2.000) -- (0.000,2.000) -- cycle;
\filldraw[cefill] (1.000,1.000) -- (2.000,1.000) -- (2.000,2.000) -- cycle;
\filldraw[cefill] (1.000,1.000) -- (2.000,2.000) -- (1.000,2.000) -- cycle;
\filldraw[cefill] (2.000,1.000) -- (3.000,1.000) -- (3.000,2.000) -- (2.000,2.000) -- cycle;
\filldraw[cefill] (3.000,1.000) -- (4.000,1.000) -- (3.000,2.000) -- cycle;
\filldraw[cefill] (4.000,1.000) -- (4.000,2.000) -- (3.000,2.000) -- cycle;
\filldraw[cefill] (4.000,1.000) -- (5.000,1.000) -- (4.000,2.000) -- cycle;
\filldraw[cefill] (5.000,1.000) -- (5.000,2.000) -- (4.000,2.000) -- cycle;
\filldraw[cefill] (5.000,1.000) -- (6.000,1.000) -- (5.000,2.000) -- cycle;
\filldraw[cefill] (6.000,1.000) -- (6.000,2.000) -- (5.000,2.000) -- cycle;
\filldraw[cefill] (0.000,2.000) -- (1.000,2.000) -- (1.000,3.000) -- (0.000,3.000) -- cycle;
\filldraw[cefill] (1.000,2.000) -- (2.000,2.000) -- (1.000,3.000) -- cycle;
\filldraw[cefill] (2.000,2.000) -- (2.000,3.000) -- (1.000,3.000) -- cycle;
\filldraw[cefill] (2.000,2.000) -- (3.000,2.000) -- (3.000,3.000) -- (2.000,3.000) -- cycle;
\filldraw[cefill] (3.000,2.000) -- (4.000,2.000) -- (5.000,2.000) -- (5.000,3.000) -- (4.000,3.000) -- (3.000,3.000) -- cycle;
\filldraw[cefill] (5.000,2.000) -- (6.000,2.000) -- (6.000,3.000) -- (5.000,3.000) -- cycle;
\filldraw[cefill] (0.000,3.000) -- (1.000,3.000) -- (1.000,4.000) -- cycle;
\filldraw[cefill] (0.000,3.000) -- (1.000,4.000) -- (0.000,4.000) -- cycle;
\filldraw[cefill] (1.000,3.000) -- (2.000,3.000) -- (3.000,3.000) -- (3.000,4.000) -- (2.000,4.000) -- (1.000,4.000) -- cycle;
\filldraw[cefill] (3.000,3.000) -- (4.000,3.000) -- (4.000,4.000) -- (3.000,4.000) -- cycle;
\filldraw[cefill] (4.000,3.000) -- (5.000,3.000) -- (4.000,4.000) -- cycle;
\filldraw[cefill] (5.000,3.000) -- (5.000,4.000) -- (4.000,4.000) -- cycle;
\filldraw[cefill] (5.000,3.000) -- (6.000,3.000) -- (6.000,4.000) -- (5.000,4.000) -- cycle;
\draw[seledge] (3.000,2.000) -- (3.000,3.000);
\draw[seledge] (2.000,2.000) -- (2.000,3.000);
\draw[seledge] (2.000,3.000) -- (3.000,3.000);
\draw[seledge] (3.000,2.000) -- (4.000,2.000);
\draw[seledge] (1.000,1.000) -- (2.000,1.000);
\draw[seledge] (1.000,2.000) -- (2.000,2.000);
\draw[seledge] (0.000,1.000) -- (1.000,1.000);
\draw[seledge] (2.000,0.000) -- (3.000,1.000);
\draw[seledge] (4.000,3.000) -- (4.000,4.000);
\draw[seledge] (5.000,2.000) -- (6.000,2.000);
\draw[seledge] (3.000,1.000) -- (4.000,1.000);
\node[vreg, dot] at (0.000,0.000) {};
\node[vreg, dot] at (1.000,0.000) {};
\node[vreg, dot] at (2.000,0.000) {};
\node[vreg, dot] at (3.000,0.000) {};
\node[vreg, dot] at (4.000,0.000) {};
\node[vreg, dot] at (5.000,0.000) {};
\node[vreg, dot] at (6.000,0.000) {};
\node[vreg, dot] at (0.000,1.000) {};
\node[vreg, dot] at (1.000,1.000) {};
\node[vreg, dot] at (2.000,1.000) {};
\node[vreg, dot] at (3.000,1.000) {};
\node[vreg, dot] at (4.000,1.000) {};
\node[vreg, dot] at (5.000,1.000) {};
\node[vreg, dot] at (6.000,1.000) {};
\node[vreg, dot] at (0.000,2.000) {};
\node[vreg, dot] at (1.000,2.000) {};
\node[vreg, dot] at (2.000,2.000) {};
\node[vreg, dot] at (3.000,2.000) {};
\node[vreg, dot] at (4.000,2.000) {};
\node[vreg, dot] at (5.000,2.000) {};
\node[vreg, dot] at (6.000,2.000) {};
\node[vreg, dot] at (0.000,3.000) {};
\node[vreg, dot] at (1.000,3.000) {};
\node[vreg, dot] at (2.000,3.000) {};
\node[vreg, dot] at (3.000,3.000) {};
\node[vreg, dot] at (4.000,3.000) {};
\node[vreg, dot] at (5.000,3.000) {};
\node[vreg, dot] at (6.000,3.000) {};
\node[vreg, dot] at (0.000,4.000) {};
\node[vreg, dot] at (1.000,4.000) {};
\node[vreg, dot] at (2.000,4.000) {};
\node[vreg, dot] at (3.000,4.000) {};
\node[vreg, dot] at (4.000,4.000) {};
\node[vreg, dot] at (5.000,4.000) {};
\node[vreg, dot] at (6.000,4.000) {};

\end{tikzpicture}
}\caption{poly}\end{subfigure}
\caption{Um exemplo de cada tipo de malha gerada pelo \code{stress.py} (grade de $6\times4$ células), com uma seleção
de arestas em vermelho. As malhas e as seleções do teste usam o mesmo gerador, com tamanhos e sementes variados.}
\label{fig:estresse-malhas}
\end{figure}

\subsection{Geração das figuras (\code{gen.py})}

\textbf{Uso:} \code{gen.py <driver> <diretório-das-figuras> <diretório-de-trabalho>}. Para cada figura, o script monta
uma malha pequena (ou lê uma do repositório), roda o \code{driver}, carrega o estado desejado e o converte num
\code{tikzpicture}, gravado em \code{figs/<nome>.tex}. Assim, toda figura de malha mostra o que o código realmente
produz; as figuras não devem ser editadas à mão, e sim regeneradas. Três detalhes do desenho:
\begin{itemize}
\item \textbf{Vista explodida.} As cópias de um vértice coincidem por definição (o código lhes dá as coordenadas do
original). Para a topologia ficar visível, cada cópia de um vértice dividido é desenhada um pouco afastada da posição
original, na direção média dos centróides dos elementos contínuos que a usam; como cada cópia é usada pelos elementos
de um setor, as cópias se afastam, cada uma para o seu setor.
\item \textbf{Arestas colapsadas.} Os dois cantos de um elemento de interface sobre uma aresta colapsada são afastados na
direção do elemento vizinho, e a aresta é desenhada como um pequeno segmento vermelho com espessura apenas visual.
\item \textbf{Buracos.} As half-edges de borda formam laços; um laço percorrido no sentido horário (área com sinal
negativa) é um buraco e é hachurado.
\end{itemize}
Os nomes dos vértices seguem a regra das figuras: o vértice original mantém o nome e suas cópias recebem índices
($K$, $K_1$, $K_2$, \dots) na ordem de criação. As figuras desenhadas à mão (gabarito, árvore de decisão, estados,
hipóteses $x_1$/$x_3$) ficam no próprio texto, em TikZ, e usam os estilos de \code{styles.tex}.
